\documentclass[../../../main.tex]{subfiles}

\begin{document}

\section{Proving and Problem Solving Tips}

Throughout studying and solving olympiad style problems, I noticed
and came up with some patterns and tips that I think could be
helpful. I am not a olympiad coach or anything, but maybe I could
share some of my thoughts and topics that I heard that quite
surprised me.

\subsection{Proving Inequalities}

Inequality is definitely one of my favorite section, and I think
there are many components that makes a solution elegant. Particularly,
I think solutions that makes inequality feel natural are very
elegant. Here is what I mean. Instead of trying to ``solve the
inequality'' by making the best out of the conditions, I think
``proving the inequality'' in a natural way is more elegant.
Say we have a constant on the right-hand side and we want to prove it.
Some ``bad proof'' could include brute forcing and trying to figure
out the meaning of the inequality just by looking at the constant.

It is, however, very important to deduce hints from the constants
because they could reveal something very important like 2025 KMO
Problem 6, which is included in Section \ref{MATH0101.26005}.
However that problem is also elegant in a way that it naturally uses
the condition to prove the inequality without brute forcing.
In short, I think solutions are elegant when we could deduce like
this is true and that is true, then this gives the constant.
Here is an example to manifest my point.

\begin{problem}{1997 KMO Problem 4 \easy}{1997 KMO Problem 4}
Given positive $a, b, c$ such that $a + b + c = abc$, show that
$\sum_\text{cyc} \frac{1}{\sqrt{1 + a^2}} \leq \frac{3}{2}$.
\end{problem}
\begin{solution}
I particularly like this problem because there are few ways of
solving this, but the standard and famous solution deduces the
constant from the left-hand side.

Consider the following equation for $\triangle{ABC}$.
\[
\tan C
= \tan{180^\circ - (A + B)}
= -\tan(A + B)
= - \frac{\tan A + \tan B}{1 - \tan A \tan B}
\]
In other words, $\tan A + \tan B + \tan C = \tan A \tan B \tan C$.
Without loss of generality, let $0^\circ < A, B < 90^\circ$.
Notice that because $c$ depends on $a$ and $b$ and $\tan A$ and
$\tan B$ spans positive reals, we can let $a = \tan A$, $b = \tan B$,
and $c = \tan C$. Moreover, we could notice that $C \in
(0^\circ, 90^\circ)$ since $c > 0$.

Continuing, consider the following inequality by Jensen.
\[
\sum_\text{cyc} \frac{1}{\sqrt{1 + \tan^2 A}}
= 3 \sum_\text{cyc} \frac{\cos A}{3}
\leq 3 \cos \left( \frac{A + B + C}{3} \right)
= \frac{3}{2}
\]
Therefore, the inequality holds.
\end{solution}

I think what makes this problem particularly elegant is the fact
that we could transform the given conditions to different forms
and we could use that to find the constant. Meaning, we could have
derived this problem from scratch without even looking at the
right-hand side.

\subsection{Solving FE and Geometry Problems}

Guess, draw well, then prove


\subsection{Proving Techniques}

Induction
    Base case, use base to go from a to a+1


\end{document}


